\chapter{Time Series Basics}
\label{ch:timeseries}

\noindent\textsc{Time series data is data with order.} Economic indicators by month, stock prices by day, temperature by year. The order matters because consecutive observations are usually correlated: today's stock price is closer to yesterday's than to a random price drawn from the historical distribution. This dependence violates the OLS independence assumption and changes almost everything about how you analyze the data.

This chapter introduces the basic vocabulary and the simplest models. Real time-series analysis is a graduate course; this is the runway.

\section{What makes time series different}

In cross-sectional data (one row per individual), shuffling the rows changes nothing. In time series, shuffling destroys the data. The order \emph{is} the information.

The defining feature is \emph{autocorrelation}: $\Cov(Y_t, Y_{t-1})$ is generally not zero. Today depends on yesterday. The deeper the dependence, the more the standard OLS toolkit fails to apply.

Three central concepts organize everything that follows:

\paragraph{Stationarity.} A series is (weakly) stationary if its mean and variance don't change over time and its autocorrelations depend only on the lag, not on $t$. Stationary series fluctuate around a constant level. Most macroeconomic series in raw form are not stationary; they trend.

\paragraph{Autocorrelation.} The correlation between $Y_t$ and $Y_{t-k}$, called the lag-$k$ autocorrelation:
\[
  \rho_k = \frac{\Cov(Y_t, Y_{t-k})}{\Var(Y_t)}.
\]
The plot of $\rho_k$ against $k$ is the \emph{autocorrelation function} (ACF). For white noise, all $\rho_k$ for $k \geq 1$ are zero (in expectation). For a random walk, $\rho_k$ decays slowly.

\paragraph{Differencing.} If $Y_t$ is non-stationary, $\Delta Y_t = Y_t - Y_{t-1}$ often is. Differencing removes trends. The number of times you have to difference a series to make it stationary is its \emph{order of integration}, written $I(d)$.

\section{The AR(1) model}

The simplest time-series model:
\begin{equation}
  Y_t = \phi Y_{t-1} + \varepsilon_t, \quad \varepsilon_t \overset{\text{iid}}{\sim} N(0, \sigma^2).
  \label{eq:ar1}
\end{equation}
Today's value is $\phi$ times yesterday's plus a fresh shock. This is called \emph{autoregressive of order 1}, AR(1).

\paragraph{Stationarity.} The AR(1) is stationary if and only if $|\phi| < 1$. The series fluctuates around zero, and its variance is finite: $\Var(Y_t) = \sigma^2/(1 - \phi^2)$. The autocorrelation function decays geometrically: $\rho_k = \phi^k$.

\paragraph{Random walk.} When $\phi = 1$, the AR(1) becomes
\[
  Y_t = Y_{t-1} + \varepsilon_t,
\]
which is the random walk. The series is non-stationary: its variance grows with $t$ ($\Var(Y_t) = t \sigma^2$ if $Y_0 = 0$). The first difference $\Delta Y_t = \varepsilon_t$ is white noise. So a random walk is $I(1)$.

\paragraph{Explosive.} When $|\phi| > 1$, the series explodes: shocks compound rather than decay. Real economic series rarely look like this for long; mean reversion or feedback eventually kicks in.

\section{Why this matters: spurious regression}

Run a regression of one random walk on another, independent random walk. The truth is no relationship. OLS will, with high probability, return a large slope and a small $p$-value anyway. This is the \emph{spurious regression} problem, and it broke decades of empirical macroeconomics before economists figured out what was happening.

The fix has two parts:
\begin{enumerate}[noitemsep]
\item Test for stationarity (Augmented Dickey-Fuller test, KPSS test).
\item If both series are non-stationary, either work in differences or use cointegration methods.
\end{enumerate}

In differences, the spurious correlation disappears, because $\Delta Y_t$ for an independent random walk is just white noise.

\section{Moving average and ARMA models}

The MA(1) model:
\[
  Y_t = \varepsilon_t + \theta \varepsilon_{t-1}.
\]
The current observation is the current shock plus a fraction of the previous shock. ACF is zero beyond lag 1.

ARMA combines AR and MA components:
\[
  Y_t = \phi_1 Y_{t-1} + \cdots + \phi_p Y_{t-p} + \varepsilon_t + \theta_1 \varepsilon_{t-1} + \cdots + \theta_q \varepsilon_{t-q}.
\]
This is ARMA($p, q$). ARIMA($p, d, q$) is the same model applied to the $d$-th difference of $Y$. The Box-Jenkins methodology is a workflow for identifying $(p, d, q)$ from data: examine the ACF and partial ACF (PACF), guess a model, check residuals, iterate.

\section{Trend, seasonality, and cycles}

Real time series are usually a mix of:
\begin{itemize}[noitemsep]
\item \textbf{Trend}: long-run direction (population growing, prices rising)
\item \textbf{Seasonality}: regular cycles tied to calendar (holiday retail, agricultural output)
\item \textbf{Cyclical fluctuations}: irregular but persistent (business cycles)
\item \textbf{Noise}: idiosyncratic shocks
\end{itemize}
Decomposing a series into these components (e.g., STL decomposition) is a standard preprocessing step. Once decomposed, you can model each component separately.

\section{Forecasting}

The point of time series modeling is often forecasting: $\hat Y_{t+h \mid t}$. For an AR(1):
\[
  \hat Y_{t+1 \mid t} = \hat\phi \cdot Y_t.
\]
Two-step ahead:
\[
  \hat Y_{t+2 \mid t} = \hat\phi \cdot \hat Y_{t+1\mid t} = \hat\phi^2 \cdot Y_t.
\]
And so on. As $h$ grows, the forecast decays toward zero (the unconditional mean), and the forecast variance grows. Long-run forecasts have wide intervals.

\begin{keyinsight}
Time series differs from cross-section because consecutive observations are typically correlated.

A series is stationary if its statistical properties don't change over time. Most economic data must be differenced or detrended before it becomes stationary. Running regressions on non-stationary series produces spurious results.

The simplest model is AR(1): today equals $\phi$ times yesterday plus noise. Stationary when $|\phi| < 1$, random walk when $\phi = 1$, explosive when $|\phi| > 1$.
\end{keyinsight}

\begin{codelab}
\begin{lstlisting}[language=Rlang]
# Simulate AR(1) processes
set.seed(2024)
n <- 200
phi_values <- c(0.3, 0.9, 1.0, -0.7)

par(mfrow = c(2, 2))
for (phi in phi_values) {
  y <- numeric(n)
  for (t in 2:n) y[t] <- phi * y[t-1] + rnorm(1)
  plot(y, type = "l", main = paste("AR(1), phi =", phi))
}

# ACF and PACF
y <- arima.sim(list(ar = 0.7), n = 300)
par(mfrow = c(1, 2))
acf(y); pacf(y)

# Fit ARIMA
fit <- arima(y, order = c(1, 0, 0))
fit
forecast::forecast(fit, h = 20)   # forecasts with intervals
\end{lstlisting}
\end{codelab}

\section{Practice problems}

\begin{enumerate}
  \item Simulate an AR(1) with $\phi = 0.5$ and one with $\phi = 0.95$. Plot both. Which one looks more like ``noise''? Which looks more like a trend? They are the same kind of model.
  \item For an AR(1) with $\phi = 0.8$ and $\sigma^2 = 1$, find $\Var(Y_t)$ in steady state and the lag-3 autocorrelation.
  \item Generate two independent random walks of length 200. Regress one on the other. Repeat 100 times and look at the distribution of $\hat\beta$ and the $p$-values. Why are most of them ``significant''?
  \item Take the differences of a random walk. Plot the differenced series and its ACF. What do you see, and why?
  \item Decompose a real time series (try \texttt{co2} or \texttt{AirPassengers} in R) using \texttt{stl()}. Identify the trend, seasonal, and remainder components.
  \item Fit AR(1), AR(2), and AR(3) models to a simulated AR(2) series. Compare AIC values. Does the lowest AIC correctly identify the true model?
\end{enumerate}

\begin{reflect}
$\triangleright$~What does it mean for a time series to be stationary? Why do we care?

\smallskip
$\triangleright$~The AR(1) with $\phi = 1$ is the random walk. Why does this single value of $\phi$ change everything about the model's behavior?

\smallskip
$\triangleright$~Why is regressing one non-stationary series on another dangerous? What does ``spurious regression'' mean?
\end{reflect}
